% Filling in this bibliographic information facilitates the
% processing of this document.
% Insert linebreaks if necessary
% The abstract page is not generated by default. The reason for this is that it
% will be generated by the unit for Publishing and Graphic sevices from the metadata 
% template ("Spikningsmallen"). If you want to include the abstract page generated from this template,
% please edit the \frontmatterMonograph or \frontmatterCS command at the en of this file
%. Because of this, some of the bibliografic data on this page will not be used.

% Abstract and titelpage
\newcommand{\authorSurname}{Maia} % Your surname
\newcommand{\authorFirstName}{Filipe} % Your given name
\newcommand{\authorFirstInitial}{RNC} % Initial of given name
\newcommand{\authorEmail}{filipe@xray.bmc.uu.se} % Your e-mail address
\newcommand{\dissertationTitle}{Ultrafast Coherent X-ray Diffractive Nano-Imaging} % The title of the dissertation
\newcommand{\dissertationSubtitle}{}% The subtitle of the dissertation (if there is any).
\newcommand{\yearOfPublication}{2010} % Year of publication
\newcommand{\placeOfDisputation}{B41, BMC, Husargatan 3, Uppsala} % Place of disputation
\newcommand{\dateOfDisputation}{Friday, May 14, 2010} % Date of disputation (Day, Month, Year)
\newcommand{\timeOfDisputation}{13:15} % Time of disputation (Swedish time format)
\newcommand{\disputationLanguage}{English} % Language the examination will be conducted in.
\newcommand{\numberOfPages}{51} % The page number of the last page
\newcommand{\placeOfPublication}{Uppsala} % The place of publication
\newcommand{\ISBN}{} % The ISBN number of the dissertation.
%\newcommand{\series}{Uppsala Dissertations from the Faculty of Science and
%  Technology} % The title of the series
\newcommand{\series}{Digital Comprehensive Summaries of Uppsala Dissertations from the Faculty of Science and Technology} % The title of the series
\newcommand{\serialNumber}{} % The number in the series.
\newcommand{\keywords}{free-electron laser, X-ray diffraction, Coherent
  Diffractive Imaging, Image Reconstruction, Single Particle Imaging, keyword6,\\ keyword7, keyword8, keyword9, keyword10,...} % Key words separated by comma
\newcommand{\department}{Department of Cell and Molecular Biology} % Name of your department
\newcommand{\departmentaddress}{Box 596, Uppsala University, SE-75124 Uppsala,
Sweden} % Address of your department
\newcommand{\ISSN}{} % (ISSN for Digital comprehensive summaries of Uppsala dissertations from the faculty of science and technology)
\newcommand{\urn}{urn:nbn:se:uu:diva-} % URN number

% List of papers
\newcommand{\listofpapers}
{\cleardoublepage
\pdfbookmark[0]{List of Publications}{LOP}
\chapter*{List of Publications}
\noindent 
\noindent This thesis is based on the following papers, which are referred
to in the text by their Roman numerals.
\vspace{1\baselineskip}

        % It is possible to refer to the published papers using labels in the text.
        % Suggested order
        % Author 1 surname, Author 2 first name initial., Author 1 surname, Author 2 first name
        % initial. etc. (Year of publication) Paper main title.
        % Paper subtitle. Name of journal in italics, volume(number):page rage
        % Example
        
    \begin{romanlist}


      \item
        \label{cowboys}
        Chapman, H.N., Barty, A., Bogan, M.J., Boutet, S., Frank, M., Hau-Riege, S.P.,
        Marchesini, S., Woods, B.W., Bajt, S., Benner, H., London, R.A., Plonjes, E., Kuhlmann,
        M., Treusch, R., Dusterer, S., Tschentscher, T., Schneider, J.R.,
        Spiller, E., Moller, T., Bostedt, C., Hoener, M., Shapiro, D.A.,
        Hodgson, K.O., Van der Spoel, D., Burmeister, F., Bergh, M., Caleman,
        C., Huldt, G., Seibert, M.M., Maia, F.R.N.C., Lee, R.W., Szoke, A.,
        Timneanu, N., Hajdu J. (2006) Femtosecond diffractive imaging with a
        soft-X-ray free-electron laser. {\em Nature Physics}, 2(12):839-843

    \item
      \label{melting}
      van der Spoel, D., Maia, F.R.N.C., Caleman, C. (2008), Structural studies of
      melting on the picosecond time scale. {\em Physical Chemistry Chemical
        Physics},10(42):6344-6349

    \item 
      \label{music}
      Ravasio, A., Gauthier, D., Maia, F.R.N.C., Billon, M., Caumes, J.P.,
      Garzella, D., Geleoc, M., Gobert, O., Hergott, J.F., Pena, A.M., Perez,
      H., Carre, B., Bourhis, E., Gierak, J., Madouri, A., Mailly, D., Schiedt,
      B., Fajardo, M., Gautier, J., Zeitoun, P., Bucksbaum, P.H., Hajdu, J.,
      Merdji, H. (2009) Single-Shot Diffractive Imaging with a Table-Top
      Femtosecond Soft X-Ray Laser-Harmonics Source. {\em Physical Review
        Letters}, 103:028104

    \item 
      \label{dynamics}
      Maia, F.R.N.C., Ekeberg, T., T\^{i}mneanu, N., van der Spoel, D., Hajdu,
      J. (2009) Structural variability and the incoherent addition of scattered
      intensities in single-particle diffraction. {\em Physical Review E}, 80:031905
      
    \item 
      \label{hawk}
      Maia, F.R.N.C.,  Ekeberg, T., van der Spoel, D., Hajdu, J. (2010)
      Hawk: the image reconstruction package for coherent X-ray diffractive
      imaging. {\em Submitted for publication}
\end{romanlist}

\vspace{1\baselineskip} \noindent {\small Reprints were made with permission from
the publishers.}

\clearpage 

\noindent {\LARGE Supporting Publications}

\vspace{1cm}

\begin{romanlist}
  \setcounter{romanlistc}{5}
\item Bergh, M., Huldt, G., Timneanu, N., Maia, F.R.N.C., Hajdu, J. (2008)
  Feasibility of imaging living cells at subnanometer resolutions by ultrafast
  X-ray diffraction. {\em Quarterly Reviews of Biophysics}, 41(3-4):181-204
  \label{QRB}
\item Bogan, M.J., Benner, W.H., Boutet, S., Rohner, U., Frank, M., Barty, A.,
  Seibert, M.M., Maia, F., Marchesini, S., Bajt, S., Woods, B., Riot, V.,
  Hau-Riege, S.P., Svenda, M., Marklund, E., Spiller, E., Hajdu, J., Chapman,
  H.N. (2008) Single particle X-ray diffractive imaging. {\em Nano Letters},
  8(1):310-316
  \label{nano}
\end{romanlist}
}


% Dedication. Dedication is optional. You can enter whatever feels right and use the fonts and images you like as long as they are embedded in the document. The suggested format does not have to be used.
\newcommand{\dedication}%
{\cleardoublepage
\thispagestyle{empty}
\vspace*{\stretch{3}}
\begin{flushright}
		
{\fontfamily{pzc}\Large\selectfont\emph{my dedication...}}

\end{flushright}
\vspace*{\stretch{1}}} % 

% Frontmatter commands
% Halftitle page for monographs
\newcommand{\halftitlepage}{%
\thispagestyle{empty}
\begin{center}
	ACTA UNIVERSITATIS UPSALIENSIS

	\emph{\series{}}
	
	\serialNumber
\end{center}
\cleardoublepage}

% Title page for monographs
\newcommand{\thetitlepage}%
{\thispagestyle{empty}
\vspace*{40mm}
\begin{center}
\Large \authorFirstName{} \authorSurname{}

\vspace{6mm}

\Huge{ \dissertationTitle{}}

\vspace{6mm}

\fontsize{14}{16}\fontshape{it}\selectfont\dissertationSubtitle{}			
\end{center}
\begin{figure}[b]
\begin{center}
\ifpdf
\includegraphics{UU_logo_pc_sv_42}
\else
\includegraphics{UU_logo_pc_sv_42.eps}
\fi
\end{center}
\end{figure}}

% Abstract. This is optional. 
\newcommand{\thesisabstract}{%
\begin{abstract}{%
	\newpage
	\thispagestyle{empty}
    \fontsize{9}{11}\selectfont
    \noindent
    \begin{flushleft}{\nohyphens
    Dissertation at Uppsala University to be publicly examined in \placeOfDisputation{}, \dateOfDisputation{} at \timeOfDisputation{}
    for the Degree of Doctor of Philosophy. The examination will be conducted in \disputationLanguage{}}\end{flushleft}

    \noindent\textbf{Abstract}
    \vspace{-9pt}
    \begin{flushleft}{\nohyphens\noindent\authorSurname{}, \authorFirstInitial{}. \yearOfPublication{}.
    \dissertationTitle{}. \dissertationSubtitle{}. Acta Universitatis
    Upsaliensis. \emph{\series{}} \serialNumber{}. \numberOfPages{}~pp.
    \placeOfPublication{}. ISBN~\ISBN{}}\end{flushleft}

% Type your abstracttext here. Remove the latin text. Try to not type more than 300 words.
\noindent Qui terempri isuam liam incum aris forteat strudentre quid inculoc iviliaet; nihilius. Catimor demquam tam, senterf rmactum ompec orum, ne culvitimis se constusquam pratastiam porum rem, que poptert batia vidiem, sendum patus? Paturo vis esum inessus, furnimi iliciem ent vidit rebaturei potiu quam ductame teris, conductam me consum retis consuniu ma, speriam rem firisum publi, quam factod norehem et grat. Maesse, se orteliciemus cem.
Fuidena ilicien ervid C. M. C. me quid rem moente, conimus nit practanti, tabena, Ti. Habunterfer que avoltum quam hinaterniam abem nonside labistiae, intelic vidi ium erce enatici ecturox nulviribem hostio nesimoeritia rei inatis.

Vivehem uoditiem imissili se in te cae ta nihili, que ina, moveres erionihil hi, Casta qui publica ercerissent, tabunici scivenatem inati, pra nest inatus habis obus cris pribus Mul vium horterioc re conduci nterum opublius bonsiliciam esus Catri sulis cons hum nius M. Viviric temnostam deatum acri cavoculatra, noveret renica; nescid C. Marit, dem ne porum pubit, nocus. Sere inatuam et, quonsi pecurori pret, senatatquo convolibus, vat, quondum prestrica amdi, co nitam ina, ad rei pri ina, num ores halarei cavenarbem.
Fuidepe fenate inimilistret patiliis facchicaet; nihilne con depos, sidemorum pro, Palicaequem pervit; ne nox mo viverfentiam re atin travo, Catquam condum oc, nit.
Tum sestiu ciemum antius concemo aedius, num il condum inatort ricaedeme fesce te int. At ore ponsula nestoratuus in popubli moveropubit; hoca; hactum perem, quem factum unte autesse ala ia ad ademquodius conducero nontis iam nes hoctum tum iam sulartest vis. Gra ni im quasta rebatui fatia? Nihilium, quitil utem et; notie ad perbis is. Satrae intrae omne imilicas oraccia me tebuncus, uterces uludem publis lia consula der liurena esimove cessupi orenamdiu moverfe terisse tured perbemne ca depons At caecermacte, vivescemur. O taberus hocus consum noxim deorbisquam quam aciis.

    \vspace{11pt}

    \noindent
    \emph{Keywords: }\keywords

    \vspace{11pt}

    \noindent
    \emph{\authorFirstName{} \authorSurname{}, \department{}, Uppsala University,
    \departmentaddress{}}

    \vspace{11pt}

    \noindent \copyright{} \authorFirstName{} \authorSurname{}
    \yearOfPublication{}

    \vspace{11pt}
    
    \noindent ISSN \ISSN{}

    \noindent ISBN \ISBN{}

    \noindent \urn{} (http://urn.kb.se/resolve?urn=\urn)}\end{abstract}}

% Title page dummy for Compehensive summaries
\newcommand{\titlepagedummy}%
{\newpage
\newpage\thispagestyle{empty}
{\noindent\LARGE Titlepage}\vspace{84pt}

\noindent This is the title page dummy. This page should be substituted for a real title page. The real title page will be provided by Publishing and Graphic services after having submitted your posting details in DiVA. The DiVA registration form can be found at \\ \href{https://uu.diva-portal.org/dream/}{https://uu.diva-portal.org/dream/}.

\vspace{1em}

\noindent More information about the publishing routines can be found at \\
\href{http://beta.ub.uu.se/pgs}{http://beta.ub.uu.se/pgs}}

%Abstractdummy
\newcommand{\abstractdummy}%
{\newpage\thispagestyle{empty}
{\noindent\LARGE Abstract page}\vspace{84pt}

\noindent This is the abstract dummy. This page should be substituted for real abstract/imprint page. The real abstract page will be provided by Publishing and Graphic services after having submitted your posting details in DiVA. The DiVA registration form can be found at \\ \href{https://uu.diva-portal.org/dream/}{https://uu.diva-portal.org/dream/}.
\vspace{1em}

\noindent More information about the publishing routines can be found at \\
\href{http://beta.ub.uu.se/pgs}{http://beta.ub.uu.se/pgs}}

% Front matter commands
\newcommand{\frontmatterMonograph}%
{\halftitlepage\thetitlepage\abstractdummy\dedication}

\newcommand{\frontmatterCS}%
{\titlepagedummy\abstractdummy\dedication\listofpapers}